\documentclass[11pt,a4paper]{article}
\usepackage[T1]{fontenc}
\usepackage{lmodern,amsmath,amssymb,amsthm,mathtools,booktabs,microtype}
\usepackage[margin=26mm]{geometry}
\usepackage[colorlinks=true,linkcolor=blue,citecolor=blue,urlcolor=blue]{hyperref}
\setlength{\emergencystretch}{2em}
\newtheorem{theorem}{Theorem}[section]
\newtheorem{proposition}[theorem]{Proposition}
\newtheorem{lemma}[theorem]{Lemma}
\newtheorem{corollary}[theorem]{Corollary}
\theoremstyle{definition}\newtheorem{definition}[theorem]{Definition}
\theoremstyle{remark}\newtheorem{remark}[theorem]{Remark}
\newcommand{\Q}{\mathbb Q}\newcommand{\Z}{\mathbb Z}
\newcommand{\R}{\mathbb R}\newcommand{\A}{\mathbb A}
\newcommand{\PP}{\mathbb P}\newcommand{\sqcl}[1]{[#1]}
\DeclareMathOperator{\Tr}{Tr}\DeclareMathOperator{\Nm}{N}
\DeclareMathOperator{\N}{N}\DeclareMathOperator{\rad}{rad}
\DeclareMathOperator{\sfpart}{sf}
\title{Limitations of Auxiliary Hypotheses in Diophantine Constructions}
\author{}\date{16 November 2026}
\begin{document}\maketitle
\begin{abstract}
Consecutive composite integers can have product with odd squarefree
part equal to two, even when they admit distinct prime divisors.
Thus distinct prime assignments alone do not control the primes
occurring to odd order. A dimension argument shows that a polynomial
injection on rational points does not furnish a finite algebraic
inverse on an open subset of affine space of dimension greater than
one. Finally, Mazur's connected-component conjecture, if true, excludes
an existential definition of $\Z$ in $\Q$. Each conclusion concerns
its stated hypothesis or encoding procedure; none is a general
variable lower bound.
\end{abstract}

\section{Scope and notation}
For a nonzero integer $n$, let
\[
                 \sfpart(n)=\prod_{p:\,v_p(n)\text{ odd}}p;
\]
the sign of $n$ is omitted. A polynomial $F\in\Z[X,\mathbf Y]$
defines $\Z$ in $\Q$ when its rational input projection is exactly
$\Z$. The results below concern three specified proposals: obtaining
odd-prime height information from distinct prime assignments, replacing
several free coordinates by a finite algebraic encoding of a polynomial
injection, and adjoining Mazur's connected-component conjecture to
premises that imply such an existential definition. The first two
limitations are unconditional; the last implication explicitly assumes
Mazur's conjecture. The height implications of $abc$, Szpiro, and
truncated Vojta are treated in~\cite{HeightHypotheses}.

\section{Prime assignments and odd valuations}

Grimm's conjecture asserts that a finite block of consecutive
composite positive integers admits distinct prime divisors,
one for each member. Its formulation and its relation to prime
gaps are discussed in~\cite{ErdosPomerance}.
The next elementary example isolates a difficulty in using
such assignments as a height estimate.

\begin{proposition}\label{prop:pell}
There are arbitrarily large positive integers $n$ such that
$n$ and $n+1$ are both composite and
\[
                    \sfpart(n(n+1))=2.
\]
Each such pair admits distinct prime divisors.
\end{proposition}
\begin{proof}
Define positive integers $a_k,b_k$ by
\[
 (17+6\sqrt8)^k=a_k+b_k\sqrt8,\qquad k\ge1.
\]
Their norm identity is $a_k^2-8b_k^2=1$, and $a_k$ is odd.
For $n_k=(a_k-1)/2$ this gives
$n_k(n_k+1)=2b_k^2$.
Multiplication by $17+6\sqrt8$ acts by the integral matrix
\[
 \begin{pmatrix}17&48\\6&17\end{pmatrix},
\]
whose determinant is one. Its image in
$\mathrm{GL}_2(\Z/144\Z)$ has some finite order $L$.
For $k\equiv1\pmod L$ we therefore have
$a_k\equiv17\pmod{144}$, whence $n_k\equiv8\pmod{72}$.
In particular $8\mid n_k$ and $9\mid n_k+1$; both are
composite. These $n_k$ tend to infinity. Their consecutive
members are coprime, so any prime divisor of one differs
from every prime divisor of the other.
\end{proof}

Proposition~\ref{prop:pell} concerns the information supplied
by a prime assignment. It is not an independence theorem
between conjectures. In particular, it shows why assigning
different primes to factors cannot be substituted for a
lower bound on primes occurring to odd order.

\section{Polynomial injections and finite algebraic encodings}

The Bombieri--Lang conjecture predicts that the rational points
of a projective variety of general type are not Zariski dense.
For surfaces, Poonen proves that this conjecture implies the
existence of a polynomial inducing an injection
$\Q^2\to\Q$~\cite[Theorem 1.1]{Poonen}.
His Remark 1.7 records a different sufficient hypothesis:
the four-term generalized $abc$ conjecture.
It is essential to distinguish that conjecture from ordinary,
three-term $abc$.

\begin{proposition}\label{prop:dimension}
Let $k$ be a characteristic-zero field and let $U$ be a nonempty
open subset of $\A^r_k$, with $r>1$.
No rational function $j:U\dashrightarrow\A^1_k$ is generically
finite onto its image.
\end{proposition}
\begin{proof}
If $j$ is constant, its only fiber has dimension $r$.
Otherwise $k(j)$ has transcendence degree one over $k$, whereas
$k(U)=k(t_1,\ldots,t_r)$ has transcendence degree $r$.
Thus
$\operatorname{trdeg}_{k(j)}k(U)=r-1>0$,
which is incompatible with a finite function-field extension.
\end{proof}

A polynomial may nevertheless be injective on $k$-rational
points, since geometric and rational fibers are different.
If $j$ is injective on $\Q^r$ and a relation is
$F(x,\mathbf y)=0$, encoding its rational witnesses by
$z=j(\mathbf y)$ gives the exact reformulation
\begin{equation}\label{eq:encoding}
 \exists\mathbf y\in\Q^r:
   F(x,\mathbf y)=0,\quad z=j(\mathbf y).
\end{equation}
Injectivity proves uniqueness of a rational inverse when one
exists; it does not supply its existence or remove the
quantifiers in~\eqref{eq:encoding}.
Proposition~\ref{prop:dimension} also excludes replacing that
inverse on an open set by a finite algebraic extension, the
setting required by the one-coordinate separation theorem.
This is a limitation of that proposed encoding. It does not
rule out a different use of polynomial injections in a
Diophantine construction.

\section{A conjecture incompatible with the projection}

Mazur's connected-component conjecture asserts that for every
variety $V/\Q$, the closure of $V(\Q)$ in $V(\mathbb R)$ has
finitely many connected components. The implication below is
classical; see~\cite[\S\S 2--3]{CornelissenZahidi}.

\begin{proposition}\label{prop:mazur}
If Mazur's connected-component conjecture holds, $\Z$ is
not Diophantine in $\Q$.
\end{proposition}
\begin{proof}
Suppose a polynomial $F(x,\mathbf y)$ defines $\Z$, and let
$V$ be its affine zero set. Let $C$ be the real topological
closure of $V(\Q)$ in $V(\mathbb R)$. The coordinate projection
$\pi:V(\mathbb R)\to\mathbb R$ is continuous. Since
$\pi(V(\Q))=\Z$ and $\Z$ is closed,
$\pi(C)\subseteq\Z$. Conversely every integer occurs at a
rational point, so $\pi(C)=\Z$.
Every connected component of $C$ has a connected image in
the discrete set $\Z$, hence a one-point image.
Finitely many components would give a finite image, a
contradiction. If the conjecture is formulated only for
irreducible varieties, apply it to the finitely many
irreducible components of $V$.
\end{proof}

Thus this topological conjecture cannot be adjoined to
arithmetic premises already implying an existential definition
of $\Z$ as a consistent method of reducing the number of witnesses.
It is distinct from the Bombieri--Lang statement used above.

\begin{thebibliography}{99}
\bibitem{ErdosPomerance}
P.~Erd\H{o}s and C.~Pomerance,
An analogue of Grimm's problem of finding distinct prime
factors of consecutive integers,
\emph{Utilitas Math.} \textbf{24} (1983), 45--65.
\url{https://math.dartmouth.edu/~carlp/PDF/38.pdf}.

\bibitem{Poonen}
B.~Poonen, Multivariable polynomial injections on rational numbers,
\emph{Acta Arith.} \textbf{145} (2010), 123--127.
\url{https://math.mit.edu/~poonen/papers/QQ.pdf}.

\bibitem{CornelissenZahidi}
G.~Cornelissen and K.~Zahidi,
Topology of Diophantine sets: remarks on Mazur's conjectures,
in \emph{Hilbert's Tenth Problem: Relations with Arithmetic
and Algebraic Geometry}, Contemp. Math.~270, 2000, 253--260.
\url{https://arxiv.org/abs/math/0006140}.

\bibitem{HeightHypotheses}
\emph{Additional arithmetic hypotheses and polynomial definitions of the integers},
project manuscript, 9 November 2026,
\nolinkurl{papers/arithmetic_hypotheses/arithmetic_hypotheses_20261109_v16.tex}.
\end{thebibliography}
\end{document}
